\section{Nesting of regional PEPS ground spaces}%
\label{sec:peps_region_ground_space_nesting}

The regional contraction with arbitrary boundary conditions of
\cite[Section~IV.C.1]{Cirac2021Matrix}, local source lines 2003--2011,
implies a compatibility between nested regions. No injectivity or
nonzero bond dimensions are needed.

\begin{theorem}[Separation of nested incident labels]%
    \label{thm:peps_nested_incident_labels}
    \lean{TNLean.PEPS.isRegionIncidentEdge_of_subset,
        TNLean.PEPS.RegionRemainderConfig,
        TNLean.PEPS.regionIncidentConfigEquivNested,
        TNLean.PEPS.sum_mul_openRegionWeight}
    \leanok
    \uses{def:peps_region_ground_space}
    Let $R\subseteq S$, and write $E_R,E_S$ for their incident
    edges. Then $E_R\subseteq E_S$, and assignments on $E_S$
    split uniquely into assignments on $E_R$ and $E_S\setminus E_R$.
    For any coefficient $f$ on the crossing-bond labels of $R$,
    contraction against its open tensor satisfies
    \begin{align}
        \sum_{\mu_R} f(\mu_R)\psi_{R,\mu_R}(\sigma_R)
         & =\sum_{\eta_R}f(\eta_R|_{\partial R})
        \prod_{v\in R} A_v(\eta_R,\sigma_R(v)).
        \label{eq:peps_nested_boundary_sum}
    \end{align}
\end{theorem}

\begin{proof}\leanok
    Restriction gives the decomposition of assignments; its inverse
    joins the two assignments on their disjoint edge sets.
    Expand the open tensor in
    (\ref{eq:peps_nested_boundary_sum}), interchange the finite
    sums, and retain the unique boundary label
    $\mu_R=\eta_R|_{\partial R}$.
\end{proof}

\begin{definition}[Physical slice of a nested region]%
    \label{def:peps_subregion_slice}
    \lean{TNLean.PEPS.subregionSliceMap,
        TNLean.PEPS.regionComplementConfigOfSubset,
        TNLean.PEPS.regionSliceMap_eq_subregionSliceMap_comp}
    \leanok
    \uses{def:peps_region_parent_hamiltonian}
    For $R\subseteq S$ and fixed physical labels $\tau$ on $R^c$,
    let $Q_{R,S,\tau}$ send a vector $\phi_S$ on $S$ to the vector
    $\sigma_R\mapsto\phi_S(\sigma_R,\tau|_{S\setminus R})$.
    Fixing all complementary labels of $R$ factors as fixing the
    labels on $S^c$ and then applying $Q_{R,S,\tau}$.
\end{definition}

\begin{theorem}[Nesting of regional ground spaces]%
    \label{thm:peps_region_ground_space_nested}
    \lean{TNLean.PEPS.openRegionWeight_nested_slice_mem_regionGroundSpace,
        TNLean.PEPS.regionGroundSpace_nested}
    \leanok
    \uses{def:peps_region_ground_space,def:peps_subregion_slice}
    If $R\subseteq S$, then every $\phi_S\in\mathcal G_S$ and
    every complementary physical configuration $\tau$ satisfy
    $Q_{R,S,\tau}\phi_S\in\mathcal G_R$.
\end{theorem}

\begin{proof}\leanok
    \uses{thm:peps_nested_incident_labels}
    It suffices to consider a boundary-labelled tensor
    $\psi_{S,\mu_S}$. Separate the incident labels into
    $\eta_R$ on $E_R$ and $\zeta$ on $E_S\setminus E_R$.
    Tensor factors at vertices in $S\setminus R$ and the boundary
    constraint on $S$ depend on $\eta_R$ only through its
    crossing labels $\mu_R$. Equation
    (\ref{eq:peps_nested_boundary_sum}) therefore gives
    \begin{align}
        Q_{R,S,\tau}\psi_{S,\mu_S}
         & =\sum_{\zeta,\mu_R}
        c_{\mu_S,\zeta,\mu_R}(\tau)\psi_{R,\mu_R},
        \label{eq:peps_nested_contraction}
    \end{align}
    where each coefficient is the product of the tensor factors
    on $S\setminus R$ when the assembled boundary labels agree
    with $\mu_S$, and is zero otherwise. The right side belongs
    to $\mathcal G_R$. Linear combinations give the assertion
    for every $\phi_S\in\mathcal G_S$.
\end{proof}

\begin{theorem}[Monotonicity of common regional conditions]%
    \label{thm:peps_parent_regions_monotonicity}
    \lean{TNLean.PEPS.regionParentGroundSpace_le_of_regions_covered,
        TNLean.PEPS.regionParentGroundSpace_antitone_regions,
        TNLean.PEPS.ker_regionParentHamiltonian_le_of_regions_covered}
    \leanok
    \uses{def:peps_region_parent_hamiltonian}
    Let $(R_i)_{i\in I}$ and $(S_j)_{j\in J}$ be collections of
    regions such that each $R_i$ is contained in some $S_j$.
    Then the common regional ground space for $(S_j)$ is contained
    in that for $(R_i)$. In particular, for finite collections,
    any positive parent Hamiltonians built from these regions
    satisfy $\ker H_S\subseteq\ker H_R$.
\end{theorem}

\begin{proof}\leanok
    \uses{thm:peps_region_ground_space_nested,
        def:peps_subregion_slice,thm:peps_region_parent_common_kernel}
    For each $R_i$, choose a containing $S_j$. A vector satisfying
    the $S_j$ condition has every complementary slice in
    $\mathcal G_{S_j}$. Apply the nesting theorem and the
    factorization of physical slices to obtain the $R_i$
    condition. The parent-kernel statement follows from the
    characterization by common regional conditions.
\end{proof}
